\documentclass[11pt]{article}
\usepackage{amsmath,amssymb,amsthm}

\newtheorem{theorem}{Theorem}

\begin{document}

\begin{theorem}[Max-Flow Min-Cut Theorem]
\label{thm:maxflowmincut}
Let $G=(V,E)$ be a directed network with a nonnegative capacity function
$u:E\to\mathbb{R}_{\ge 0}$, together with a source $s$ and a sink $t$. The
value of an $s$--$t$ flow is the total amount of flow leaving $s$. An
$s$--$t$ cut is a subset $S\subseteq V$ such that $s\in S$ and $t\notin S$,
and its capacity is
\[
\operatorname{cap}(S)=\sum_{(i,j)\in E,\ i\in S,\ j\notin S} u(i,j).
\]
Then
\[
\max\{\text{value of a feasible } s\text{--}t \text{ flow}\}
=
\min\{\operatorname{cap}(S): s\in S,\ t\notin S\}.
\]
In other words, the maximum value of a feasible flow equals the minimum
capacity of an $s$--$t$ cut.
\end{theorem}

\begin{proof}
Let $x$ be any feasible $s$--$t$ flow, and let its value be $z$. For any
$s$--$t$ cut $S$, sum the flow-conservation equations over all vertices in
$S\setminus\{s\}$. All contributions on edges with both endpoints in $S$
cancel, and we obtain the cut identity
\[
z
=
\sum_{(i,j)\in E,\ i\in S,\ j\notin S} x_{ij}
-
\sum_{(j,i)\in E,\ j\notin S,\ i\in S} x_{ji}.
\]
Since $0\le x_{ij}\le u(i,j)$ on every edge, it follows that
\[
z\le
\sum_{(i,j)\in E,\ i\in S,\ j\notin S} u(i,j)
=
\operatorname{cap}(S).
\]
Thus every feasible flow has value at most the capacity of every $s$--$t$ cut.

For the reverse inequality, run the augmenting-path algorithm until it stops,
and let $x^*$ be the resulting flow. In the residual network of $x^*$, let
$S$ be the set of vertices reachable from $s$. Because there is no augmenting
path, $t\notin S$, so $S$ is an $s$--$t$ cut.

Now consider an edge $(i,j)$ with $i\in S$ and $j\notin S$. If it were not
saturated, then the residual network would contain the forward edge
$i\to j$, so $j$ would also be reachable from $s$, a contradiction. Hence
$x^*_{ij}=u(i,j)$ for every such edge.

Next consider an edge $(j,i)$ with $j\notin S$ and $i\in S$. If
$x^*_{ji}>0$, then the residual network would contain the reverse edge
$i\to j$, again making $j$ reachable from $s$, a contradiction. Hence
$x^*_{ji}=0$ for every edge entering $S$ from $V\setminus S$.

Substituting these two facts into the cut identity gives
\[
\operatorname{value}(x^*)
=
\sum_{(i,j)\in E,\ i\in S,\ j\notin S} u(i,j)
=
\operatorname{cap}(S).
\]
So there exists a feasible flow whose value equals the capacity of some
$s$--$t$ cut. Combined with the first part, this proves that the maximum flow
value equals the minimum cut capacity.
\end{proof}

\end{document}
